\documentclass[../main.tex]{subfiles}
\begin{document}

\section{Probleme}

\subsection{Problema Bart (NerdArena)}
\label{problem:bart}

\href{https://www.nerdarena.ro/problema/bart}{enunț}
$\bullet$
\hyperref[code:bart]{sursă}

Problema ne cere să aflăm cel mai mic prefix al unui șir $S$ pe care, dacă îl repetăm, obținem șirul $S$. Ultima repetiție poate fi incompletă.

Fie răspunsul $k$. Atunci $S[0 \dots k-1] = S[k \dots 2k-1] = \dots$. Dar aceasta înseamnă că $S[0 \dots n - k - 1] = S[k \dots n - 1]$. Cu alte cuvinte, sufixul de lungime $n - k$ este și prefix! Răspunsul este pur și simplu $n - \pi[n - 1]$.

\subsection{Problema Cifru (Baraj ONI 2006)}
\label{problem:cifru}

\href{https://kilonova.ro/problems/1862}{enunț}
$\bullet$
\hyperref[code:cifru]{sursă}

În această problemă recunoaștem testul clasic dacă două șiruri $S$ și $T$ se pot obține unul dintr-altul prin rotații. Dar cum gestionăm permutarea?

Vom modifica algoritmul KMP ca să țină cont și de permutări. Să considerăm șablonul din exemplu, $P = (2\ 1\ 3\ 3\ 2\ 1)$. La poziția $i$ dorim să aflăm lungimea celui mai lung prefix care \textit{este} și sufix. Dar aceasta depinde de sensul cuvîntului „este”\footnote{Exact ca în cazul lui Bill Clinton.}. Nu ne interesează ca prefixul să fie \textbf{egal} cu sufixul, ci să existe o permutare care mapează sufixul în prefix. De exemplu, $\pi[4]=2$ pentru că prefixul $(2\ 1)$ se potrivește cu sufixul $(3\ 2)$.

Cum construim funcția $\pi$ astfel modificată? Să pornim de la $\pi[2] = 2$ (deoarece $(2\ 1)$ se potrivește cu $(1\ 3)$). Am putea extinde cu succes perechea de subșiruri cu o poziție dacă $(2\ 1\ 3)$ s-ar potrivi cu $(1\ 3\ 3)$. Dar nu este cazul. În primul subșir, 3 este o valoare nou întîlnită, pe cînd în al doilea valoarea 3 a apărut și pe poziția 1. De aceea nu există nicio permutare care să mapeze prefixul peste sufix. Așadar, trebuie să înlocuim testul de egalitate de caractere din algoritmul KMP normal cu un test de asemănare între caractere:

\begin{itemize}
  \item Dacă $P[k]$ este la prima apariție (deci $P[k] \neq P[0], P[1], ..., P[k-1]$), trebuie ca și $P[i] \neq P[i - k], P[i - k + 1], ..., P[i - 1]$.
  \item Altfel, dacă $P[k] = P[k']$ pentru un $k' < k$, trebuie ca și $P[i] = P[i - (k - k')]$.
\end{itemize}

Putem verifica aceste (in)egalități dacă stocăm, pentru fiecare poziție $k$, poziția apariției anterioare a lui $P[k]$ în $P$.

Același raționament îl facem și în a doua trecere, în care îl căutăm pe $P$ în $S + S$. Cînd evaluăm $S[i]$ și sîntem în starea $k$, vom confrunta aparițiile anterioare ale lui $S[i]$ în $S$ cu aparițiile lui $P[k]$ în $P$.

\subsection{Problema Anthem of Berland (Codeforces)}
\label{problem:anthem-of-berland}

\href{https://kilonova.ro/problems/1862}{enunț}
$\bullet$
\hyperref[code:anthem-of-berland]{sursă}

Problema ne poate duce rapid cu gîndul la elementele de bază:

\begin{itemize}
  \item KMP și funcția prefix, deoarece avem un șablon ale cărui apariții se pot suprapune.
  \item Programare dinamică.
\end{itemize}

Să punem la punct detaliile. Voi folosi simbolul $P$ (\textit{pattern}) în loc de $T$, din obișnuință. Ne putem convinge rapid că o abordare greedy nu funcționează. Iată un exemplu similar cu cel din enunț: $S$ = \texttt{????abcab}, $P$ = \texttt{abcab}. Atunci este important să completăm șirul ca $S$ = \texttt{?abcabcab}, începînd cu a doua poziție, pentru a obține două potriviri.

Ce parametri ne trebuie ca să definim bine recurența? Din exemplul de mai sus observăm că nu este suficient să definim $i$ = numărul de litere consumate din $S$ ($0 \leq i \leq |S|$). Trebuie să putem spune „după prima literă, vrem să fim în starea 0, și abia de la a doua începînd să trecem în stările 1, 2 etc.” Deci mai definim și $j$ = starea în care ne aflăm ($0 \leq j \leq |P|$). Fie $d_{i,j}$ numărul maxim de suprapuneri pe care îl putem obține pînă la aceste coordonate sau -1 dacă nu putem ajunge la aceste coordonate.

Eu am ales să formulez programarea dinamică înainte, astfel: din $d_{i,j}$ trecem în $d_{i+1,k+1}$, unde:

\begin{itemize}
  \item $k$ este orice tranziție posibilă din starea $j$;

  \item $S[i] = P[k]$ sau $S[i] = \texttt{?}$;

  \item Cînd spunem „trecem în”, vrem să spunem că $d_{i+1,k+1}$ primește maximul din ce avea deja și $d_{i,j}$, eventual plus încă o unitate dacă $k +1 = |P|$ (deoarece am completat încă o apariție).
\end{itemize}

De exemplu, pentru șirul $P$ = \texttt{abacaba}, valoarea $d_{20}{7}$ arată cîte potriviri putem obține în primele 20 de caractere, astfel încît după al 20-lea (deci $S[19]$) să fim în starea 7 (deci să terminăm o potrivire). De aici putem trece:

\begin{itemize}
  \item în $d_{21,4}$ dacă următorul caracter este \texttt{c} sau \texttt{?}, deoarece din \texttt{abacabac} putem refolosi 4 litere;

  \item în $d_{21,2}$ dacă următorul caracter este \texttt{b} sau \texttt{?}, deoarece din \texttt{abacabab} putem refolosi 2 litere;

  \item în $d_{21,1}$ dacă următorul caracter este \texttt{a} sau \texttt{?}, deoarece din \texttt{abacabaa} putem refolosi o literă;

  \item în $d_{21,0}$, altfel.
\end{itemize}

\subsubsection*{Soluție în \texorpdfstring{$\bigoh(SP^2)$}{O(SP2)}}

Prima iterație a codului meu proceda chiar astfel. Pentru fiecare $i$ și fiecare $j$, pornea cu $k = j$ și repeta $k = \pi[k - 1]$ pînă cînd ajungea la 0. Dar acest cod este prea lent, deoarece ajunge la $\bigoh(SP^2)$ atunci cînd funcția $\pi$ scade lent, de exemplu pentru un șablon ca $P$ = \texttt{aaa \dots aaa}.

\subsubsection*{Soluție în \texorpdfstring{$\bigoh(SP\Sigma)$}{SP Sigma}}

În loc de aceasta, putem să calculăm tranzițiile din starea $d_{i,j}$ în alte stări conform următorului caracter, care fie este impus (literă), fie poate lua $\Sigma$ valori (\texttt{?}). Practic, trebuie să răspundem la întrebări de tipul: „Dacă ne aflăm în starea $j$, iar următorul caracter de la intrare este $S[i]$, în ce stare ajungem?”. Aceste informații reprezintă exact matricea de tranziție despre care am vorbit în capitolul despre automate! Cu ajutorul funcției de prefix $\pi$, putem calcula această matrice în $\bigoh(P^2)$ (\href{https://codeforces.com/contest/808/submission/376033109}{o sursă} mai veche) sau chiar $\bigoh(P \Sigma)$.

\end{document}
